% Telos shared preamble.
%
% Every Telos publication is designed to be printed on a home laser printer,
% one sheet at a time, and carried into a boat or a garage. Two constraints
% follow from that and govern everything below:
%
%   1. Pure monochrome. No value is ever a hue. Emphasis is carried by rule
%      weight, fill, and type, never by color, so a grayscale or black-only
%      printer loses nothing.
%   2. Card density. A "one-pager" is a hard contract: the leaf must fit on a
%      single side of US Letter. The layout is therefore tight by default and
%      the environments below are sized to leave room for a plate.

\documentclass[10pt,letterpaper]{article}

\usepackage[margin=0.55in,includefoot,footskip=0.24in]{geometry}
\usepackage[T1]{fontenc}
\usepackage[utf8]{inputenc}
\usepackage{lmodern}
\usepackage{microtype}
\usepackage{array}
\usepackage{booktabs}
\usepackage{longtable}
\usepackage{tabularx}
\usepackage{enumitem}
\usepackage{needspace}
\usepackage{multicol}
\usepackage{ragged2e}
\usepackage{graphicx}
\usepackage{wrapfig}
\usepackage{xcolor}
\usepackage{fancyhdr}
\usepackage{titlesec}
\usepackage{hyperref}
\usepackage[most]{tcolorbox}
\usepackage{tikz}
\usetikzlibrary{arrows.meta,calc,positioning,shapes.geometric,decorations.pathmorphing,patterns,fadings}

% Semantic names are kept so a document reads intelligibly, but every value is
% black, white, or a neutral gray that survives a black-only print driver.
\definecolor{telosink}{gray}{0}
\definecolor{telospaper}{gray}{1}
\definecolor{telosrule}{gray}{0.35}
\definecolor{telosfill}{gray}{0.90}
\definecolor{telosdeep}{gray}{0.72}
\definecolor{telosmute}{gray}{0.45}

\hypersetup{hidelinks,pdfborder={0 0 0}}

% Keep an unchanged source byte-reproducible across rebuilds: no build-clock
% creation date and no random trailer ID.
\pdfinfoomitdate=1
\pdftrailerid{}

\setlength{\parindent}{0pt}
\setlength{\parskip}{0.42em}
\setlength{\tabcolsep}{4pt}
\setlength{\columnsep}{1.1em}
\renewcommand{\arraystretch}{1.18}
\setlist[itemize]{leftmargin=1.15em,itemsep=0.12em,topsep=0.2em,parsep=0pt}
\setlist[enumerate]{leftmargin=1.5em,itemsep=0.18em,topsep=0.2em,parsep=0pt}
\setlist[description]{leftmargin=0pt,itemsep=0.16em,topsep=0.2em,parsep=0pt,
  font=\normalfont\bfseries}

% ---------------------------------------------------------------------------
% Document identity. Each leaf declares these before \begin{document} so the
% running head, the colophon, and the PDF metadata all agree.
% ---------------------------------------------------------------------------
\newcommand{\TelosProject}{Telos}
\newcommand{\TelosDocTitle}{}
\newcommand{\TelosDocSubtitle}{}
\newcommand{\TelosRevision}{}

\newcommand{\telosmeta}[4]{%
  \renewcommand{\TelosProject}{#1}%
  \renewcommand{\TelosDocTitle}{#2}%
  \renewcommand{\TelosDocSubtitle}{#3}%
  \renewcommand{\TelosRevision}{#4}%
  \hypersetup{pdftitle={#2},pdfsubject={#3; version #4},
    pdfkeywords={Telos, version #4},pdfauthor={Telos}}%
}

\pagestyle{fancy}
\fancyhf{}
\renewcommand{\headrulewidth}{0pt}
\renewcommand{\footrulewidth}{0.4pt}
\fancyfoot[L]{\footnotesize\textsc{\TelosProject}}
\fancyfoot[C]{\footnotesize\TelosDocTitle}
\fancyfoot[R]{\footnotesize\thepage}

% The first page of a one-pager carries the masthead instead of a running foot,
% so its footer is the compact provenance line.
\fancypagestyle{telosfirst}{%
  \fancyhf{}%
  \renewcommand{\headrulewidth}{0pt}%
  \renewcommand{\footrulewidth}{0.4pt}%
  \fancyfoot[L]{\footnotesize\textsc{\TelosProject}}%
  \fancyfoot[C]{\footnotesize\TelosRevision}%
  \fancyfoot[R]{\footnotesize\thepage}%
}

% ---------------------------------------------------------------------------
% Masthead. A heavy rule, the title, a light subtitle, a thin rule. The
% optional argument is a right-aligned tag (a species' scientific name, a
% lesson number) that sits on the title's baseline.
% ---------------------------------------------------------------------------
\newcommand{\telosmasthead}[1][]{%
  \thispagestyle{telosfirst}%
  \begingroup
  \setlength{\parskip}{0pt}%
  \noindent\rule{\linewidth}{1.6pt}\par
  \vspace{0.30em}
  \noindent
  \begin{minipage}[b]{0.66\linewidth}
    \raggedright{\LARGE\bfseries\TelosDocTitle}
  \end{minipage}%
  \hfill
  \begin{minipage}[b]{0.32\linewidth}
    \raggedleft{\small\itshape #1}
  \end{minipage}\par
  \vspace{0.18em}
  \noindent{\small\TelosDocSubtitle}\par
  \vspace{0.28em}
  \noindent\rule{\linewidth}{0.4pt}\par
  \endgroup
  \vspace{0.35em}
}

% ---------------------------------------------------------------------------
% Headings. Compact, small-caps, rule-anchored — a card has no room for the
% article class's default vertical air.
% ---------------------------------------------------------------------------
\titleformat{\section}{\normalfont\large\bfseries}{}{0pt}{}[\vspace{-0.55em}\rule{\linewidth}{0.4pt}]
\titlespacing*{\section}{0pt}{0.85em}{0.35em}
\titleformat{\subsection}{\normalfont\normalsize\bfseries}{}{0pt}{}
\titlespacing*{\subsection}{0pt}{0.6em}{0.2em}
\titleformat{\subsubsection}{\normalfont\small\bfseries\scshape}{}{0pt}{}
\titlespacing*{\subsubsection}{0pt}{0.45em}{0.12em}

% A short run-in heading for card interiors that must not consume a full line.
\newcommand{\telosrubric}[1]{\textbf{#1}\enspace}

% Cross-reference helper. Every Telos project is a set of loose printed sheets,
% so a reference names the sheet rather than a page number.
\newcommand{\sheetref}[1]{\textit{see the \textbf{#1} sheet}}

% Which decision, ADR or sheet governs the paragraph you are reading. A reader
% who disagrees with an instruction should be able to find the reasoning behind
% it rather than having to take it on trust.
\newcommand{\governedby}[1]{{\footnotesize\textsc{governed by #1}}}

% ---------------------------------------------------------------------------
% Quantities. siunitx is not in this TeX Live split, and the guides need only a
% fixed thin space and upright units, so define that directly rather than take
% a package dependency the documented install would not satisfy.
%   \qty{9}{kV}  ->  9 kV        \unitof{kV}  ->  kV
% ---------------------------------------------------------------------------
\newcommand{\unitof}[1]{\ensuremath{\mathrm{#1}}}
\newcommand{\qty}[2]{#1\,\unitof{#2}}
\newcommand{\numrange}[2]{#1\,--\,#2}
\newcommand{\qtyrange}[3]{#1\,--\,#2\,\unitof{#3}}
% Inches and feet appear constantly in the fishing guide; keep one spelling.
\newcommand{\inch}[1]{#1\,in}
\newcommand{\feet}[1]{#1\,ft}

% ---------------------------------------------------------------------------
% Cards and boxes.
% ---------------------------------------------------------------------------

% Titled card with a hairline frame. The default card for facts a reader looks
% up rather than reads through.
\newtcolorbox{teloscard}[1]{
  enhanced, breakable,
  colback=telospaper, colframe=telosink, colbacktitle=telospaper,
  coltitle=telosink,
  boxrule=0.5pt, sharp corners,
  left=6pt, right=6pt, top=4pt, bottom=4pt,
  toptitle=3pt, bottomtitle=3pt, titlerule=0.4pt,
  before skip=6pt, after skip=6pt,
  title={#1}, fonttitle=\small\bfseries
}

% Untitled left-ruled block: hierarchy without a redundant label.
\newtcolorbox{telosframe}{
  enhanced, breakable,
  colback=telospaper, frame hidden, boxrule=0pt, sharp corners,
  borderline west={1.3pt}{0pt}{telosink},
  left=8pt, right=2pt, top=3pt, bottom=3pt,
  before skip=6pt, after skip=6pt
}

% Filled emphasis panel — the "read this first" block. The 10% fill stays
% legible under a cheap laser toner curve and still reads as distinct.
\newtcolorbox{telospanel}[1]{
  enhanced, breakable,
  colback=telosfill, colframe=telosink, colbacktitle=telosfill,
  coltitle=telosink,
  boxrule=0.5pt, sharp corners,
  left=6pt, right=6pt, top=4pt, bottom=4pt,
  toptitle=3pt, bottomtitle=3pt, titlerule=0.4pt,
  before skip=6pt, after skip=6pt,
  title={#1}, fonttitle=\small\bfseries
}

% Hazard block. Deliberately the heaviest object on any page: a 1.4pt frame
% plus a solid title bar. Nothing else in Telos is allowed to look like this,
% so a reader skimming a lesson cannot mistake it for ordinary emphasis.
\newtcolorbox{teloshazard}[1]{
  enhanced, breakable,
  colback=telospaper, colframe=telosink,
  colbacktitle=telosink, coltitle=telospaper,
  boxrule=1.4pt, sharp corners,
  left=6pt, right=6pt, top=5pt, bottom=5pt,
  toptitle=3pt, bottomtitle=3pt,
  before skip=8pt, after skip=8pt,
  title={\MakeUppercase{#1}}, fonttitle=\small\bfseries
}

% A stop-and-check gate. Used before anything destructive, and before anything
% that would be expensive to discover later. Shared by every project that has
% a procedure someone could get hurt following carelessly.
\newtcolorbox{gate}[1]{
  enhanced, breakable,
  colback=telospaper, colframe=telosink,
  boxrule=1.0pt, sharp corners,
  left=6pt, right=6pt, top=4pt, bottom=4pt,
  toptitle=3pt, bottomtitle=3pt, titlerule=0.4pt,
  before skip=7pt, after skip=7pt,
  colbacktitle=telospaper, coltitle=telosink,
  title={\MakeUppercase{#1}}, fonttitle=\small\bfseries
}

% ---------------------------------------------------------------------------
% Rating dots. A five-step scale that prints identically in black-only: filled
% versus open circles, never a shaded bar.
% ---------------------------------------------------------------------------
\newcommand{\telosdot}{\tikz[baseline=-0.42ex]\fill[telosink] (0,0) circle (0.28ex);}
\newcommand{\telosopendot}{\tikz[baseline=-0.42ex]\draw[telosink,line width=0.28pt] (0,0) circle (0.28ex);}
\makeatletter
\newcounter{telos@rate}
\newcommand{\telosrating}[1]{%
  \begingroup
  \setcounter{telos@rate}{0}%
  \@whilenum\value{telos@rate}<#1\do{\telosdot\kern0.22em\stepcounter{telos@rate}}%
  \@whilenum\value{telos@rate}<5\do{\telosopendot\kern0.22em\stepcounter{telos@rate}}%
  \endgroup
}
\makeatother

% ---------------------------------------------------------------------------
% Tables.
%
% tabularx reads its body by scanning for a literal \end{tabularx}, so it can
% never be wrapped in a \newenvironment. Documents therefore open tabularx and
% longtable directly and take their house style from these column types and the
% booktabs rules. The one exception is the fact strip below, which is built on
% plain tabular precisely so it *can* be an environment.
% ---------------------------------------------------------------------------
\newcolumntype{L}{>{\raggedright\arraybackslash}X}
\newcolumntype{R}{>{\raggedleft\arraybackslash}X}
\newcolumntype{C}{>{\centering\arraybackslash}X}
\newcolumntype{B}{>{\bfseries\raggedright\arraybackslash}X}
\newcolumntype{N}{>{\raggedright\arraybackslash}p{0.16\linewidth}}

% Key/value fact strip. Two flush columns of short lookups; the label column is
% a fixed fraction so a stack of cards aligns across the whole guide.
\newenvironment{telosfacts}[1][0.24]
  {\begingroup\small
   \setlength{\parskip}{0pt}%
   \renewcommand{\arraystretch}{1.25}%
   \begin{tabular}{@{}>{\bfseries\raggedright\arraybackslash}p{\dimexpr#1\linewidth\relax}@{\hspace{0.7em}}>{\raggedright\arraybackslash}p{\dimexpr\linewidth-#1\linewidth-0.7em\relax}@{}}}
  {\end{tabular}\par\endgroup}
\newcommand{\telosfact}[2]{#1 & #2\\}

% ---------------------------------------------------------------------------
% Seasonal / diurnal activity bar.
%
% \telosseasonbar{4,3,2,1,1,2,3,3,4,5,4,3} draws a twelve-cell month strip
% where each value 0-5 is a fill height. Height, not gray value, carries the
% signal, so it is exact under any toner curve; the cell is outlined either way
% so an empty month is still visibly a month.
% ---------------------------------------------------------------------------
\newcommand{\telosseasonlabels}{J,F,M,A,M,J,J,A,S,O,N,D}
\newcommand{\telosseasonbar}[1]{%
  \begingroup
  \begin{tikzpicture}[x=1.25em,y=2.1ex,baseline=0pt]
    \foreach \v [count=\i from 0] in {#1} {
      \draw[telosrule,line width=0.3pt] (\i,0) rectangle (\i+0.86,5);
      \ifnum\v>0
        \fill[telosink] (\i,0) rectangle (\i+0.86,\v);
      \fi
    }
    \foreach \m [count=\i from 0] in \telosseasonlabels {
      \node[font=\tiny,anchor=north,inner sep=1.2pt] at (\i+0.43,-0.1) {\m};
    }
  \end{tikzpicture}%
  \endgroup
}

% Four-cell day-part strip: dawn, midday, dusk, night.
\newcommand{\telosdaylabels}{Dawn,Mid,Dusk,Night}
\newcommand{\telosdaybar}[1]{%
  \begingroup
  \begin{tikzpicture}[x=2.9em,y=2.1ex,baseline=0pt]
    \foreach \v [count=\i from 0] in {#1} {
      \draw[telosrule,line width=0.3pt] (\i,0) rectangle (\i+0.92,5);
      \ifnum\v>0
        \fill[telosink] (\i,0) rectangle (\i+0.92,\v);
      \fi
    }
    \foreach \m [count=\i from 0] in \telosdaylabels {
      \node[font=\tiny,anchor=north,inner sep=1.2pt] at (\i+0.46,-0.1) {\m};
    }
  \end{tikzpicture}%
  \endgroup
}

% A bar needs vertical room for its labels when it sits in a fact strip.
\newcommand{\telosbarrow}[1]{\rule{0pt}{4.2ex}#1\rule[-1.6ex]{0pt}{1.6ex}}

% ---------------------------------------------------------------------------
% Plates. A species or apparatus illustration with a caption rule beneath it.
% \telosplate{width fraction}{file}{caption}
% ---------------------------------------------------------------------------
\newcommand{\telosplate}[3]{%
  \begingroup
  \centering
  \includegraphics[width=#1\linewidth]{#2}\par
  \vspace{0.15em}
  \begin{minipage}{#1\linewidth}
    \rule{\linewidth}{0.4pt}\par
    \vspace{0.1em}
    {\footnotesize #3\par}
  \end{minipage}\par
  \endgroup
  \vspace{0.3em}
}

% TikZ drawing conventions shared by every rig, knot, and circuit diagram.
\tikzset{
  telos line/.style={draw=telosink,line width=0.7pt},
  telos hair/.style={draw=telosink,line width=0.35pt},
  telos heavy/.style={draw=telosink,line width=1.2pt},
  telos node/.style={draw=telosink,fill=telospaper,align=center,inner sep=4pt,font=\small},
  telos emphasis/.style={draw=telosink,line width=1.2pt,fill=telospaper,align=center,
    inner sep=4pt,font=\small\bfseries},
  telos label/.style={font=\scriptsize,inner sep=1.5pt},
  telos arrow/.style={-{Stealth[length=1.8mm]},draw=telosink,line width=0.6pt},
  telos dim/.style={|<->|,draw=telosink,line width=0.35pt},
  telos water/.style={draw=telosrule,line width=0.4pt,decorate,
    decoration={snake,amplitude=0.4mm,segment length=3mm}},
}

% ---------------------------------------------------------------------------
% Lesson furniture
% ---------------------------------------------------------------------------

% \lessonhead{number}{time}{who does what}
\newcommand{\lessonhead}[3]{%
  \par\noindent
  \begin{minipage}{\linewidth}
    \rule{\linewidth}{0.4pt}\par
    \vspace{0.3em}
    {\small\textbf{Lesson #1}\quad$\cdot$\quad #2\quad$\cdot$\quad #3\par}
    \vspace{0.25em}
    \rule{\linewidth}{0.4pt}
  \end{minipage}\par
  \vspace{0.5em}}

% The materials list. Two columns of short lines; a lesson that needs a long
% shopping list is a lesson that has not been thought through.
\newenvironment{materials}
  {\begin{telospanel}{What you need}\begin{multicols}{2}
   \begin{itemize}[leftmargin=1em,itemsep=0.05em,topsep=0pt]}
  {\end{itemize}\end{multicols}\end{telospanel}}

% The four rubrics every lesson carries, in the same order every time:
% what you are about to see, what to do, what it means, and what to ask.
\newcommand{\lessonsection}[1]{\section*{#1}\vspace{-0.35em}\rule{\linewidth}{0.4pt}\vspace{0.2em}}

% A prediction box. The point of the curriculum is that they commit to an
% answer before the demonstration, not after.
\newtcolorbox{predict}{
  enhanced, breakable,
  colback=telospaper, colframe=telosink,
  boxrule=0.5pt, sharp corners,
  left=6pt, right=6pt, top=4pt, bottom=4pt,
  before skip=6pt, after skip=6pt,
  borderline west={2.2pt}{0pt}{telosink},
  title={\small\bfseries Predict first}, fonttitle=\small\bfseries,
  colbacktitle=telospaper, coltitle=telosink,
  toptitle=3pt, bottomtitle=3pt, titlerule=0.4pt,
}

% ---------------------------------------------------------------------------
% Colophon. Rendered once, at the end of the last page, in the recovered footer
% area so it never creates a page of its own.
% ---------------------------------------------------------------------------
\newif\ifTelosColophonRendered
\newcommand{\TelosColophonExtra}{}
\newcommand{\TelosColophon}{%
  \ifTelosColophonRendered\else
  \global\TelosColophonRenderedtrue
  \par
  \enlargethispage{0.5in}%
  \vspace{0.3em}%
  \begingroup
  \setlength{\parskip}{0pt}%
  \begin{minipage}{\linewidth}
  \rule{\linewidth}{0.35pt}\par
  \vspace{0.3em}%
  \fontsize{7}{8.2}\selectfont
  \textbf{Telos.} \TelosProject{} — \TelosDocTitle. Revised \TelosRevision.
  \TelosColophonExtra{}
  Project-created text and diagrams are licensed \href{https://creativecommons.org/licenses/by/4.0/}{CC BY 4.0}.
  Incorporated public-domain artwork remains public domain and is credited on
  the plate it appears with. See \texttt{LICENSE} and \texttt{CREDITS.md} in the
  source. Nothing here is professional advice; verify current regulations and
  follow the stated safety procedure.
  \end{minipage}
  \endgroup
  \fi
}
\AtEndDocument{\TelosColophon}

% Shared macros for the homelab project.
%
% This project's documents are read in two very different situations: at a desk
% while designing, and standing in front of a broken machine at eleven at night
% with no network. The second case is the one that sets the style. Commands are
% set apart so they can be found by eye, every dangerous step has a visible
% verification after it, and nothing depends on being able to reach a web page.

% ---------------------------------------------------------------------------
% Decision-record entries, used by the generated digest.
% ---------------------------------------------------------------------------
\newcommand{\adrentry}[4]{%
  \par\vspace{0.7em}
  \needspace{4\baselineskip}%
  \noindent
  \begin{minipage}{\linewidth}
    \rule{\linewidth}{0.6pt}\par
    \vspace{0.25em}
    {\small\textbf{ADR #1}\quad #2\par}
    \vspace{0.1em}
    {\fontsize{7}{8.2}\selectfont\textsc{#3}\quad$\cdot$\quad #4\par}
    \vspace{0.2em}
    \rule{\linewidth}{0.3pt}
  \end{minipage}\par
  \vspace{0.3em}}

% ---------------------------------------------------------------------------
% Procedures.
%
% A rebuild manual is a sequence of commands and the observable evidence that
% each one worked. These environments keep that pairing structural rather than
% a matter of the author remembering.
% ---------------------------------------------------------------------------

% A block of shell to type. Deliberately not a listing package: plain verbatim
% in a ruled box prints identically everywhere and has no font dependency.
\newenvironment{shellblock}
  {\par\vspace{0.25em}\noindent
   \begin{minipage}{\linewidth}
   \rule{\linewidth}{0.3pt}\par\vspace{0.15em}
   \begingroup\footnotesize\ttfamily\obeylines\obeyspaces\parindent=0pt\parskip=0pt}
  {\endgroup\par\vspace{0.15em}\rule{\linewidth}{0.3pt}\end{minipage}\par\vspace{0.25em}}

% What you should see if the previous step worked. A step without one of these
% is a step you cannot verify, and in a rebuild manual that is a defect.
\newcommand{\expectthis}[1]{%
  \par\vspace{0.1em}
  \noindent\begin{minipage}{\linewidth}
  \hspace*{0.6em}\begin{minipage}{\dimexpr\linewidth-0.6em\relax}
  {\footnotesize\textbf{Expect:} #1\par}
  \end{minipage}\end{minipage}\par\vspace{0.3em}}

% \begin{gate} now lives in common/preamble.tex, shared with the other
% projects that document procedures.

% \governedby now lives in common/preamble.tex.

% An instance value that comes from the private overlay. In the published
% documents these render as visible placeholders, which is deliberate: a
% published manual must be complete and obviously generic, never look like it
% is describing somebody's actual network.
\newcommand{\instancevalue}[2]{\texttt{\textlangle #1\textrangle}%
  \footnote{Instance value. Supplied from \texttt{homelab/instance/} at build
  time; #2.}}
\newcommand{\placeholder}[1]{\texttt{\textlangle #1\textrangle}}

% ---------------------------------------------------------------------------
% Role and boundary diagrams.
% ---------------------------------------------------------------------------
\tikzset{
  role/.style={draw=telosink,line width=0.9pt,fill=telospaper,align=center,
    inner sep=5pt,font=\small,minimum width=2.6cm},
  roleoff/.style={draw=telosmute,line width=0.5pt,dash pattern=on 2.5pt off 1.8pt,
    fill=telospaper,align=center,inner sep=5pt,font=\small,minimum width=2.6cm},
  boundary/.style={draw=telosmute,line width=0.5pt,dash pattern=on 3pt off 2pt,
    rounded corners=3pt},
  flow/.style={-{Stealth[length=2mm]},draw=telosink,line width=0.7pt},
  hlabel/.style={font=\fontsize{6.4}{7.4}\selectfont},
}


\telosmeta{Homelab}{Workstation Owner Guide}
  {Normal use, automatic updates, recovery, and useful evidence}{20260727.001}
\newcommand{\ownercheck}[2]{\par\needspace{3\baselineskip}\noindent
  \begin{tabularx}{\linewidth}{@{}p{.30\linewidth}L@{}}
  {\footnotesize\textbf{Do:} #1} & {\footnotesize\textbf{Check:} #2}\\
  \end{tabularx}\vspace{.16em}}

\begin{document}
\telosmasthead[Keep this copy on the laptop]

This is the first-stop guide for someone using a Telos dual-boot workstation
away from home. It covers ordinary use and safe, reversible recovery. It
contains no household names, addresses, credentials, recovery keys, or
administrator-only procedure.

\begin{gate}{The safe boundary}
You may restart, choose either operating system, reconnect Wi-Fi, check update
status, free your own ordinary storage, test optional storage, and collect the
evidence listed here. Stop before changing firmware, partitions, boot entries,
directory membership, system time by hand, package databases, security policy,
or another person's files.
\end{gate}

\section{Know the normal state}

\begin{center}
\begin{tikzpicture}[node distance=8mm and 9mm,font=\scriptsize]
  \node[role] (power) {power on};
  \node[role,right=of power] (menu) {five-second\\boot menu};
  \node[role,above right=5mm and 10mm of menu] (win) {Windows 11\\default};
  \node[role,below right=5mm and 10mm of menu] (arch) {Arch Linux\\choose explicitly};
  \node[role,right=15mm of win] (local) {local profile\\opens};
  \node[role,right=15mm of arch] (local2) {local home\\opens};
  \draw[flow] (power)--(menu);
  \draw[flow] (menu)--(win);
  \draw[flow] (menu)--(arch);
  \draw[flow] (win)--(local);
  \draw[flow] (arch)--(local2);
\end{tikzpicture}
\end{center}

Windows starts if nobody touches the menu. Select Arch before the countdown
ends when you need it. Your files live on the laptop, so login must work
without home Wi-Fi, the Controller, or optional network storage.

\ownercheck{Sign in once while connected, disconnect Wi-Fi, lock, and sign in
again.}{The same account opens the same local profile without a network error.
Repeat once in each operating system before a long trip.}
\ownercheck{Save work and sign out. Shut down before packing the laptop.}{No
update screen, activity light, or warm fan remains when it enters the bag.}

\begin{gate}{Offline login has a limit}
A previously used account may continue to open this laptop while disconnected,
even after an administrator disables its network access. This protects travel
use but is not remote erasure. Report a lost laptop promptly.
\end{gate}

\section{Connect to Wi-Fi}

Use the managed-workstation network supplied with the laptop. Its credential
may be unique to this machine. Do not copy it to phones, televisions, guests,
or another laptop.

\begin{enumerate}
\item Turn Wi-Fi off, wait five seconds, and turn it on.
\item Select the saved managed-workstation network.
\item Open one ordinary HTTPS site.
\item Check that the clock and time zone are correct automatically.
\item At home, sign out and in once to prove directory contact.
\end{enumerate}

\expectthis{Internet access works. Management pages, peer devices, and
unrelated household networks may remain unreachable by design.}
\failurebranch{If the network is absent, test one other known network or a
phone hotspot. If other networks work, record ``managed network absent.'' If
none work, restart once and collect the Wi-Fi evidence on page
\pageref{sec:evidence}. Do not forget and recreate the saved network unless an
administrator asks.}

\section{Automatic updates}

\textbf{Windows:} leave Windows Update enabled. Windows downloads and installs
security and quality updates automatically, then asks for or schedules a
restart. Plug in overnight at least weekly. Never use an advertised
third-party ``driver updater.''

\textbf{Arch:} a daily system timer attempts one complete
\texttt{pacman -Syu} transaction. It runs only on AC power, with at least
8 GiB free, no other package transaction, and the official mirror reachable.
A missed run is retried after the next boot. It records before/after package
lists and verifies installed package files. It never interrupts a user session
for a reboot.

\begin{center}
\begin{tikzpicture}[node distance=7mm,font=\scriptsize]
  \node[role] (timer) {daily timer};
  \node[role,right=of timer] (gate) {power + space\\network + lock};
  \node[role,right=of gate] (upgrade) {one full\\upgrade};
  \node[role,right=of upgrade] (verify) {verify +\\save evidence};
  \node[roleoff,below=of gate] (later) {defer safely\\try later};
  \draw[flow] (timer)--(gate);
  \draw[flow] (gate)--node[above]{ready}(upgrade);
  \draw[flow] (upgrade)--(verify);
  \draw[flow] (gate)--node[right]{not ready}(later);
\end{tikzpicture}
\end{center}

\ownercheck{Plug in, connect to the internet, and leave the laptop awake for
an hour each week.}{Windows shows no pending security update. On Arch,
\texttt{systemctl status homelab-arch-update.timer} says \texttt{active
(waiting)} and shows a next run.}
\ownercheck{Restart after an update requests it, or when Arch reports a reboot
recommendation. Reopen your work first.}{The selected OS boots, Wi-Fi
reconnects, the clock is correct, and login succeeds.}

\begin{gate}{Never partially update Arch}
Do not run \texttt{pacman -Sy}, force package replacement, delete the pacman
lock, change mirrors during a transaction, or power off while packages are
being installed. Arch is maintained by full-system upgrades.
\end{gate}

\section{Optional network storage}

Optional storage is a convenience, not your home directory. It appears after
login as a mapped drive in Windows or a mounted folder under your Arch home.
Login must succeed when storage is absent.

\begin{tabularx}{\linewidth}{@{}p{.24\linewidth}p{.32\linewidth}L@{}}
\toprule
\textbf{What you see} & \textbf{Likely meaning} & \textbf{Safe response}\\
\midrule
Drive/folder opens & Network and authorization work & Use normally\\
Storage unavailable & Away, server down, or route blocked & Work locally;
retry later\\
Access denied & Server reachable; identity or permission differs & Record the
exact message; do not alter permissions\\
Login itself fails & Not an acceptable storage failure & Disconnect Wi-Fi,
test cached login, escalate\\
\bottomrule
\end{tabularx}

\ownercheck{Before travel, disconnect from the home network and open ordinary
local documents.}{The desktop and local files open promptly; the optional
share may show unavailable without delaying login.}
\ownercheck{Copy important work to its intended backup location when
available.}{Open the copied file from the destination. A visible name alone
does not prove the bytes arrived.}

\section{Recovery ladder}

Use this order. Stop as soon as normal operation returns; do not stack
unrelated fixes.

\begin{enumerate}
\item \textbf{Protect the work.} Save open files locally. Photograph the exact
message if saving is impossible.
\item \textbf{Check simple state.} Connect power. Check airplane mode, Wi-Fi,
free storage, date, time zone, and Caps Lock.
\item \textbf{Retry once.} Repeat the failed action and note any change.
\item \textbf{Restart normally.} Use Restart, not a forced power-off.
\item \textbf{Try the other OS.} This separates hardware or firmware trouble
from one operating system.
\item \textbf{Try offline login.} Disconnect Wi-Fi and use an account that has
previously signed in on this machine.
\item \textbf{Use local rescue only when directed.} It regains local access;
it does not administer the directory. Never send its password.
\item \textbf{Escalate with evidence.} Stop before firmware, boot, disk,
identity, or security changes.
\end{enumerate}

\begin{tabularx}{\linewidth}{@{}p{.24\linewidth}p{.35\linewidth}L@{}}
\toprule
\textbf{Symptom} & \textbf{One safe test} & \textbf{Stop boundary}\\
\midrule
No power & Known-good outlet and charger; hold power once & Burning smell,
swelling, liquid, unusual heat\\
No boot menu & Remove docks and USB devices; power on once & Firmware settings
or boot-entry editing\\
One OS fails & Boot the other OS; record the failing stage & Repair media,
partition, or bootloader commands\\
Password rejected & Check layout/Caps Lock; test offline once & Repeated
guesses, account tools, clock edits\\
No internet & Test another known network & Network reset or driver replacement\\
Disk nearly full & Review your trash/downloads & System paths, another user's
files, partitions\\
Update failed & Leave power/network connected; record status & Forced package
removal, lock deletion, firmware update\\
\bottomrule
\end{tabularx}

\begin{gate}{Immediate stop}
Power down, unplug, and seek help for a swollen battery, liquid ingress, smoke,
burning odor, crackling, repeated thermal shutdown, exposed conductors, or a
charger/port too hot to touch. Never recharge a visibly damaged battery.
\end{gate}

\section{Collect useful evidence without secrets}
\label{sec:evidence}

Create a plain-text note named with the date and machine label. Include:

\begin{itemize}
\item OS, whether it reached the desktop, and whether power was connected;
\item local date/time shown and approximate location;
\item what you attempted, exact result, and whether it worked before;
\item exact error text or a photo cropped to the relevant window;
\item last successful boot, login, and update;
\item whether restart, the other OS, offline login, and another network worked;
\item Windows Update history or these read-only Arch results:
\end{itemize}

\begin{shellblock}
systemctl status homelab-arch-update.timer --no-pager
systemctl status homelab-arch-update.service --no-pager
journalctl -u homelab-arch-update.service -b --no-pager
df -h /
\end{shellblock}

\expectthis{The commands report timer state, the most recent update attempt,
and free space. They do not change the machine.}

\textbf{Remove before sharing:} passwords; Wi-Fi keys or QR codes; recovery
keys; tokens; private keys; browser contents; personal documents; full
addresses; public IP addresses; serial numbers; and another person's name or
files. Send the bundle only through the family's agreed private channel. Never
paste it into a public issue or chat.

\askbefore{Can another person reproduce the symptom from this note? Does it
say what succeeded as well as what failed? Did you open every file and image
to confirm no secret or unrelated personal material is visible?}

\section{When to ask for help}

\textbf{Ask soon:} update failure persists through the next powered,
internet-connected attempt; free space stays below 8 GiB; time repeatedly
drifts; optional storage says access denied; either OS stops booting; or login
works only offline.

\textbf{Ask immediately:} possible hardware damage; both operating systems
fail; the disk disappears; a recovery key is unexpectedly requested; an
unknown administrator prompt appears; the laptop is lost or stolen; or you
suspect account compromise.

\textbf{Administrator work only:} directory accounts and revocation, domain
rejoin, rescue-account rotation, firmware and Secure Boot, disk layout,
encryption, boot entries, PXE, network policy, update policy, package repair,
and storage authorization.

\section{Keep this guide available offline}

Every workstation image should install this versioned PDF at
\path{/usr/local/share/doc/telos/workstation-owner-guide.pdf} on Arch and in
the shared Telos documentation folder on Windows, with a desktop or Start-menu
shortcut. The public web copy is convenient but is not the recovery copy.

Controller convergence distributes a newer reviewed guide by version; it
never requires reinstalling the workstation. Deployment copies a new file,
verifies its SHA-256 digest, then replaces the shortcut atomically. Keep the
previous copy until the new PDF opens locally.

\ownercheck{Disconnect every network and open the shortcut in each operating
system.}{The title is ``Workstation Owner Guide,'' footer version is
\texttt{20260727.001}, all pages render, and the recovery ladder needs no link.}
\ownercheck{After an update, compare its version with the published
manifest.}{The new digest matches and the previous readable copy remains until
acceptance is complete.}

\begin{gate}{Before leaving home}
Boot both operating systems, complete automatic updates, restart each once,
test offline login, open local files without optional storage, connect to one
non-home network, and open this guide with all networks disabled.
\end{gate}

\section{Owner handoff record}

\begin{tabularx}{\linewidth}{@{}p{.38\linewidth}L@{}}
\toprule
Machine label & \rule{0pt}{2.5ex}\hrulefill\\
Guide version verified & \rule{0pt}{2.5ex}\hrulefill\\
Windows boot/update/offline login & \rule{0pt}{2.5ex}\hrulefill\\
Arch boot/update/offline login & \rule{0pt}{2.5ex}\hrulefill\\
Non-home Wi-Fi tested & \rule{0pt}{2.5ex}\hrulefill\\
Optional storage fails soft & \rule{0pt}{2.5ex}\hrulefill\\
Private help channel recorded offline & \rule{0pt}{2.5ex}\hrulefill\\
Owner and administrator review date & \rule{0pt}{2.5ex}\hrulefill\\
\bottomrule
\end{tabularx}

\end{document}
